% Options for packages loaded elsewhere
\PassOptionsToPackage{unicode}{hyperref}
\PassOptionsToPackage{hyphens}{url}
%
\documentclass[
]{article}
\usepackage{amsmath,amssymb}
\usepackage{iftex}
\ifPDFTeX
  \usepackage[T1]{fontenc}
  \usepackage[utf8]{inputenc}
  \usepackage{textcomp} % provide euro and other symbols
\else % if luatex or xetex
  \usepackage{unicode-math} % this also loads fontspec
  \defaultfontfeatures{Scale=MatchLowercase}
  \defaultfontfeatures[\rmfamily]{Ligatures=TeX,Scale=1}
\fi
\usepackage{lmodern}
\ifPDFTeX\else
  % xetex/luatex font selection
\fi
% Use upquote if available, for straight quotes in verbatim environments
\IfFileExists{upquote.sty}{\usepackage{upquote}}{}
\IfFileExists{microtype.sty}{% use microtype if available
  \usepackage[]{microtype}
  \UseMicrotypeSet[protrusion]{basicmath} % disable protrusion for tt fonts
}{}
\makeatletter
\@ifundefined{KOMAClassName}{% if non-KOMA class
  \IfFileExists{parskip.sty}{%
    \usepackage{parskip}
  }{% else
    \setlength{\parindent}{0pt}
    \setlength{\parskip}{6pt plus 2pt minus 1pt}}
}{% if KOMA class
  \KOMAoptions{parskip=half}}
\makeatother
\usepackage{xcolor}
\usepackage[margin=1in]{geometry}
\usepackage{color}
\usepackage{fancyvrb}
\newcommand{\VerbBar}{|}
\newcommand{\VERB}{\Verb[commandchars=\\\{\}]}
\DefineVerbatimEnvironment{Highlighting}{Verbatim}{commandchars=\\\{\}}
% Add ',fontsize=\small' for more characters per line
\usepackage{framed}
\definecolor{shadecolor}{RGB}{248,248,248}
\newenvironment{Shaded}{\begin{snugshade}}{\end{snugshade}}
\newcommand{\AlertTok}[1]{\textcolor[rgb]{0.94,0.16,0.16}{#1}}
\newcommand{\AnnotationTok}[1]{\textcolor[rgb]{0.56,0.35,0.01}{\textbf{\textit{#1}}}}
\newcommand{\AttributeTok}[1]{\textcolor[rgb]{0.13,0.29,0.53}{#1}}
\newcommand{\BaseNTok}[1]{\textcolor[rgb]{0.00,0.00,0.81}{#1}}
\newcommand{\BuiltInTok}[1]{#1}
\newcommand{\CharTok}[1]{\textcolor[rgb]{0.31,0.60,0.02}{#1}}
\newcommand{\CommentTok}[1]{\textcolor[rgb]{0.56,0.35,0.01}{\textit{#1}}}
\newcommand{\CommentVarTok}[1]{\textcolor[rgb]{0.56,0.35,0.01}{\textbf{\textit{#1}}}}
\newcommand{\ConstantTok}[1]{\textcolor[rgb]{0.56,0.35,0.01}{#1}}
\newcommand{\ControlFlowTok}[1]{\textcolor[rgb]{0.13,0.29,0.53}{\textbf{#1}}}
\newcommand{\DataTypeTok}[1]{\textcolor[rgb]{0.13,0.29,0.53}{#1}}
\newcommand{\DecValTok}[1]{\textcolor[rgb]{0.00,0.00,0.81}{#1}}
\newcommand{\DocumentationTok}[1]{\textcolor[rgb]{0.56,0.35,0.01}{\textbf{\textit{#1}}}}
\newcommand{\ErrorTok}[1]{\textcolor[rgb]{0.64,0.00,0.00}{\textbf{#1}}}
\newcommand{\ExtensionTok}[1]{#1}
\newcommand{\FloatTok}[1]{\textcolor[rgb]{0.00,0.00,0.81}{#1}}
\newcommand{\FunctionTok}[1]{\textcolor[rgb]{0.13,0.29,0.53}{\textbf{#1}}}
\newcommand{\ImportTok}[1]{#1}
\newcommand{\InformationTok}[1]{\textcolor[rgb]{0.56,0.35,0.01}{\textbf{\textit{#1}}}}
\newcommand{\KeywordTok}[1]{\textcolor[rgb]{0.13,0.29,0.53}{\textbf{#1}}}
\newcommand{\NormalTok}[1]{#1}
\newcommand{\OperatorTok}[1]{\textcolor[rgb]{0.81,0.36,0.00}{\textbf{#1}}}
\newcommand{\OtherTok}[1]{\textcolor[rgb]{0.56,0.35,0.01}{#1}}
\newcommand{\PreprocessorTok}[1]{\textcolor[rgb]{0.56,0.35,0.01}{\textit{#1}}}
\newcommand{\RegionMarkerTok}[1]{#1}
\newcommand{\SpecialCharTok}[1]{\textcolor[rgb]{0.81,0.36,0.00}{\textbf{#1}}}
\newcommand{\SpecialStringTok}[1]{\textcolor[rgb]{0.31,0.60,0.02}{#1}}
\newcommand{\StringTok}[1]{\textcolor[rgb]{0.31,0.60,0.02}{#1}}
\newcommand{\VariableTok}[1]{\textcolor[rgb]{0.00,0.00,0.00}{#1}}
\newcommand{\VerbatimStringTok}[1]{\textcolor[rgb]{0.31,0.60,0.02}{#1}}
\newcommand{\WarningTok}[1]{\textcolor[rgb]{0.56,0.35,0.01}{\textbf{\textit{#1}}}}
\usepackage{graphicx}
\makeatletter
\def\maxwidth{\ifdim\Gin@nat@width>\linewidth\linewidth\else\Gin@nat@width\fi}
\def\maxheight{\ifdim\Gin@nat@height>\textheight\textheight\else\Gin@nat@height\fi}
\makeatother
% Scale images if necessary, so that they will not overflow the page
% margins by default, and it is still possible to overwrite the defaults
% using explicit options in \includegraphics[width, height, ...]{}
\setkeys{Gin}{width=\maxwidth,height=\maxheight,keepaspectratio}
% Set default figure placement to htbp
\makeatletter
\def\fps@figure{htbp}
\makeatother
\setlength{\emergencystretch}{3em} % prevent overfull lines
\providecommand{\tightlist}{%
  \setlength{\itemsep}{0pt}\setlength{\parskip}{0pt}}
\setcounter{secnumdepth}{-\maxdimen} % remove section numbering
\ifLuaTeX
  \usepackage{selnolig}  % disable illegal ligatures
\fi
\usepackage{bookmark}
\IfFileExists{xurl.sty}{\usepackage{xurl}}{} % add URL line breaks if available
\urlstyle{same}
\hypersetup{
  pdftitle={MicrobTiSDA-vignette},
  hidelinks,
  pdfcreator={LaTeX via pandoc}}

\title{MicrobTiSDA-vignette}
\author{}
\date{\vspace{-2.5em}}

\begin{document}
\maketitle

\section{Introduction to MicrobTiSDA}\label{introduction-to-microbtisda}

Longitudinal analysis of microbiomes is crucial for understanding the
temporal dynamics of host-associated or environmental microbiomes,
Laying the foundation for exploring ecological patterns and mechanisms
of host-microbiome interactions. Sampling microbiota at multiple time
points allows researchesrs to uncover community trends, stability and
responses to external interventions such as therapies. Thus, developing
efficient and accurate methods for microbiome time-series analysis has
become increasingly urgent. Here, we developed MicrobTiSDA, an R package
that integrates the Learning interactions from Microbial Time-Series
(LIMITS) algorithm. By employing discrete-time Lotka-Volterra models,
MicrobTiSDA infers inter-species interactions from microbiome
time-series data. Moreover, the package enables the construction of
regression models to investigate the dynamic abundance changes of
individual species within microbial communities and provides reliable
methods for identifying species groups with similar (or opposing)
temporal pattens, thereby offering novel perspectives on potential
inter-species relationships.

\subsection{Getting started}\label{getting-started}

First we have to load the installed package and other necessary packages
in the workspace with

\begin{Shaded}
\begin{Highlighting}[]
\FunctionTok{library}\NormalTok{(}\StringTok{"MicrobTiSDA"}\NormalTok{)}
\FunctionTok{library}\NormalTok{(}\StringTok{"tidyr"}\NormalTok{)}
\FunctionTok{library}\NormalTok{(}\StringTok{"dplyr"}\NormalTok{)}
\FunctionTok{library}\NormalTok{(}\StringTok{"ggplot2"}\NormalTok{)}
\end{Highlighting}
\end{Shaded}

\subsubsection{Case study 1: Infant gut microbiome
dataset}\label{case-study-1-infant-gut-microbiome-dataset}

The first case study involves applying MicrobTiSDA to an infant gut
microbiome dataset. We focused on data from dizygotic twins (ID10 and
ID11) collected from day six to day sixty after birth. This dataset
comprises a total of 77 fecal samples with 476,921 operational taxonomic
units (OTUs). ID10 was hospitalized due to an infection with
Streptococcus serogroup B β-hemolytic (S. agalactiae) and received
antibiotic treatment from day 22 to day 33. Overall, we collected 77
fecal samples, with ID10 contributing 34 samples and ID11 providing 43
samples.

The input data required by MicrobTiSDA typically includes a species
composition data table from high-throughput sequencing (OTU/ASV table),
a metadata file that records all sample information, and species
annotation data for microbial characteristics. The input species
composition table is required to have microbial feature IDs as rows and
sample IDs as columns. The input metadata must have sample IDs as row
names and must include information such as the subject ID to which the
samples belong and the sampling time of the samples.

Here, we load the data of this case study.

\begin{Shaded}
\begin{Highlighting}[]
\FunctionTok{data}\NormalTok{(}\StringTok{"OSLO.infant.data"}\NormalTok{)}
\FunctionTok{data}\NormalTok{(}\StringTok{"OSLO.infant.meta"}\NormalTok{)}
\FunctionTok{data}\NormalTok{(}\StringTok{"OSLO.infant.taxa"}\NormalTok{)}
\end{Highlighting}
\end{Shaded}

We first filter out OTUs from the loaded OTU table that have a total
abundance of less than 5000 across all samples and a prevalence of less
than 20\% in all samples from each subject.

\begin{Shaded}
\begin{Highlighting}[]
\NormalTok{oslo\_data\_filt }\OtherTok{=} \FunctionTok{Data.filter}\NormalTok{(}\AttributeTok{Data =}\NormalTok{ OSLO.infant.data,}\AttributeTok{metadata =}\NormalTok{ OSLO.infant.meta,}
                             \AttributeTok{OTU\_counts\_filter\_value =} \DecValTok{5000}\NormalTok{, }\AttributeTok{OTU\_filter\_value =} \FloatTok{0.2}\NormalTok{,}
                             \AttributeTok{Group\_var =} \StringTok{\textquotesingle{}ID\textquotesingle{}}\NormalTok{)}
\end{Highlighting}
\end{Shaded}

Subsequently, we performe interpolation on the filtered OTU table to
generate continuous time series data for each subject from the first
sampling time point to just before the last sampling time point.

\begin{Shaded}
\begin{Highlighting}[]
\NormalTok{oslo\_data\_intp }\OtherTok{=} \FunctionTok{Data.interpolate}\NormalTok{(}\AttributeTok{Data =}\NormalTok{ oslo\_data\_filt,}\AttributeTok{metadata =}\NormalTok{ OSLO.infant.meta,}
                                  \AttributeTok{Group\_var =} \StringTok{\textquotesingle{}ID\textquotesingle{}}\NormalTok{,}
                                  \AttributeTok{Sample\_ID =} \StringTok{\textquotesingle{}ID\textquotesingle{}}\NormalTok{,}
                                  \AttributeTok{Sample\_Time =} \StringTok{\textquotesingle{}day\textquotesingle{}}\NormalTok{)}
\end{Highlighting}
\end{Shaded}

The interpolated OTU table and the new metadata will be used for
subsequent analyses.

\begin{Shaded}
\begin{Highlighting}[]
\NormalTok{oslo\_data\_int }\OtherTok{=}\NormalTok{ oslo\_data\_intp}\SpecialCharTok{$}\NormalTok{Interpolated\_Data}
\NormalTok{oslo\_meta\_int }\OtherTok{=}\NormalTok{ oslo\_data\_intp}\SpecialCharTok{$}\NormalTok{Interpolated\_Data\_metadata}
\end{Highlighting}
\end{Shaded}

Transform the filtered and interpolated data using modified centered
log-ratio transformation.

\begin{Shaded}
\begin{Highlighting}[]
\NormalTok{oslo\_data\_trans }\OtherTok{=} \FunctionTok{Data.trans}\NormalTok{(}\AttributeTok{Data =}\NormalTok{ oslo\_data\_int,}
                             \AttributeTok{metadata =}\NormalTok{ oslo\_meta\_int,}
                             \AttributeTok{Group\_var =} \StringTok{\textquotesingle{}Group\textquotesingle{}}\NormalTok{)}
\end{Highlighting}
\end{Shaded}

Next, we will infer the species interactions within the gut microbiomes
of each individual and visualize the topological results of these
species interactions. To reduce computation time, we will set the number
of iterations for bagging to 5 (the default is 10).

\begin{Shaded}
\begin{Highlighting}[]
\NormalTok{oslo\_data\_interact }\OtherTok{=} \FunctionTok{Spec.interact}\NormalTok{(}\AttributeTok{Data =}\NormalTok{ oslo\_data\_trans,}
                                   \AttributeTok{metadata =}\NormalTok{ oslo\_meta\_int,}
                                   \AttributeTok{Group\_var =} \StringTok{\textquotesingle{}Group\textquotesingle{}}\NormalTok{,}
                                   \AttributeTok{num\_iterations =} \DecValTok{5}\NormalTok{)}
\end{Highlighting}
\end{Shaded}

Visualize the topological results of interspecies interactions within
microbiomes of each individual.

\begin{Shaded}
\begin{Highlighting}[]
\NormalTok{oslo\_interact\_vis }\OtherTok{=} \FunctionTok{Interact.vis}\NormalTok{(}\AttributeTok{Interact\_data =}\NormalTok{ oslo\_data\_interact,}
                                 \AttributeTok{count\_data =}\NormalTok{ oslo\_data\_trans,}
                                 \AttributeTok{metadata =}\NormalTok{ oslo\_meta\_int,}
                                 \AttributeTok{Subject\_ID =} \StringTok{\textquotesingle{}ID\textquotesingle{}}\NormalTok{,}
                                 \AttributeTok{Interact\_threshold =} \FloatTok{1e{-}6}\NormalTok{,}
                                 \AttributeTok{legend\_text\_size =} \DecValTok{8}\NormalTok{,}
                                 \AttributeTok{legend\_title\_size =} \DecValTok{10}\NormalTok{,}
                                 \AttributeTok{title\_size =} \DecValTok{12}\NormalTok{,}
                                 \AttributeTok{label\_distance =} \FloatTok{0.35}\NormalTok{,}
                                 \AttributeTok{label\_font\_size =} \DecValTok{3}\NormalTok{)}
\NormalTok{oslo\_interact\_vis}\SpecialCharTok{$}\NormalTok{ID\_10}
\NormalTok{oslo\_interact\_vis}\SpecialCharTok{$}\NormalTok{ID\_11}
\end{Highlighting}
\end{Shaded}

Another functionality of MicrobTiSDA is the construction of natural
spline regression models for the time series data of each microbial
feature. Before this, it is necessary to construct the design matrix.
The sampling time data for each sample will be added as columns to the
species composition data frame after transformation. Additionally, dummy
variables will be used to label each subject's samples, where 1
indicates that the sample belongs to that subject and 0 indicates that
it does not.

\begin{Shaded}
\begin{Highlighting}[]
\NormalTok{oslo\_data\_design }\OtherTok{=} \FunctionTok{Design}\NormalTok{(}\AttributeTok{metadata =}\NormalTok{ oslo\_meta\_int,}
                          \AttributeTok{Group\_var =} \StringTok{\textquotesingle{}Group\textquotesingle{}}\NormalTok{,}
                          \AttributeTok{Sample\_ID =} \StringTok{\textquotesingle{}ID\textquotesingle{}}\NormalTok{,}
                          \AttributeTok{Sample\_Time =} \StringTok{\textquotesingle{}Time\textquotesingle{}}\NormalTok{,}
                          \AttributeTok{Pre\_processed\_Data =}\NormalTok{ oslo\_data\_trans)}
\end{Highlighting}
\end{Shaded}

Here, we construct natural spline regression models for each microbial
feature based on the transformed time series data. Given that subject
ID10 was hospitalized and received antibiotic treatment from day 22 to
day 33 after birth, we will use days 22 and 33 as the time konts.
Additionally, we will filter out microbial features from each subject's
gut microbiome time series that have fewer than 10 unique transformed
abundance values, as such features do not provide sufficient data for
constructing regression models.

\begin{Shaded}
\begin{Highlighting}[]
\NormalTok{oslo\_model }\OtherTok{=} \FunctionTok{Reg.SPLR}\NormalTok{(}\AttributeTok{Data\_for\_Reg =}\NormalTok{ oslo\_data\_design,}
                      \AttributeTok{pre\_processed\_data =}\NormalTok{ oslo\_data\_trans,}
                      \AttributeTok{Knots =} \FunctionTok{c}\NormalTok{(}\DecValTok{22}\NormalTok{,}\DecValTok{33}\NormalTok{),}
                      \AttributeTok{unique\_values =} \DecValTok{10}\NormalTok{)}
\end{Highlighting}
\end{Shaded}

Next, we will use the constructed regression models for the microbial
feature time series to make predictions, aiming to investigate the
temporal patterns of the microbial features.

\begin{Shaded}
\begin{Highlighting}[]
\NormalTok{oslo\_model\_pred }\OtherTok{=} \FunctionTok{Pred.data}\NormalTok{(}\AttributeTok{Fitted\_models =}\NormalTok{ oslo\_model,}
                            \AttributeTok{metadata =}\NormalTok{ oslo\_meta\_int,}
                            \AttributeTok{Group =} \StringTok{\textquotesingle{}Group\textquotesingle{}}\NormalTok{,}
                            \AttributeTok{Sample\_Time =} \StringTok{\textquotesingle{}Time\textquotesingle{}}\NormalTok{,}
                            \AttributeTok{time\_step =} \DecValTok{1}\NormalTok{)}
\end{Highlighting}
\end{Shaded}

By calculating the correlation distance between the predicted microbial
feature time series from the computational model, MicrobTiSDA employs a
hierarchical clustering method to cluster OTUs based on the correlation
distance between them.

\begin{Shaded}
\begin{Highlighting}[]
\NormalTok{oslo\_model\_clust }\OtherTok{=} \FunctionTok{Data.cluster}\NormalTok{(}\AttributeTok{predicted\_data =}\NormalTok{ oslo\_model\_pred,}
                                \AttributeTok{clust\_method =} \StringTok{\textquotesingle{}average\textquotesingle{}}\NormalTok{,}
                                \AttributeTok{dend\_title\_size =} \DecValTok{12}\NormalTok{,}
                                \AttributeTok{font\_size =} \DecValTok{2}\NormalTok{)}
\end{Highlighting}
\end{Shaded}

After obtaining the dendrogram of OTU clustering, we manually set a
correlation distance threshold of less than 0.15 (\code{cut_height} =
0.15) to group OTUs into clusters based on the output dendrogram.

\begin{Shaded}
\begin{Highlighting}[]
\NormalTok{oslo\_clust\_results }\OtherTok{=} \FunctionTok{Data.cluster.cut}\NormalTok{(}\AttributeTok{cluster\_outputs =}\NormalTok{ oslo\_model\_clust,}\AttributeTok{cut\_height =} \FloatTok{0.15}\NormalTok{,}\AttributeTok{font\_size =} \DecValTok{2}\NormalTok{)}
\end{Highlighting}
\end{Shaded}

If the \texttt{auto\_cutree} parameter of the Data.cluster function is
set to FALSE, a dendrogram of the microbial feature time series clusters
for each subject will be output during the execution of the
\texttt{Data.cluster} function.

Next, we can visualize the clustered OTU groups. The temporal patterns
of the OTU clusters for each subject will be generated separately in the
plots. If the \texttt{Taxa} parameter is set to \texttt{NULL}, the plots
will only display the IDs of each microbial feature. If \texttt{Taxa}
contains species annotation information, the annotation will be added in
parentheses following the microbial feature IDs in the plots. If
\texttt{plot\_dots} is \texttt{True}, the original transformed microbial
feature abundances will be visualized in figures.

\begin{Shaded}
\begin{Highlighting}[]
\NormalTok{oslo\_model\_vis }\OtherTok{=} \FunctionTok{Data.visual}\NormalTok{(}\AttributeTok{cluster\_results =}\NormalTok{ oslo\_clust\_results,}
                             \AttributeTok{cutree\_by =} \StringTok{\textquotesingle{}height\textquotesingle{}}\NormalTok{,}
                             \AttributeTok{cluster\_height =} \FunctionTok{c}\NormalTok{(}\FloatTok{0.15}\NormalTok{,}\FloatTok{0.15}\NormalTok{),}
                             \AttributeTok{predicted\_data =}\NormalTok{ oslo\_model\_pred,}
                             \AttributeTok{Design\_data =}\NormalTok{ oslo\_data\_design,}
                             \AttributeTok{pre\_processed\_data =}\NormalTok{ oslo\_data\_trans,}
                             \AttributeTok{plot\_dots =} \ConstantTok{TRUE}\NormalTok{,}
                             \AttributeTok{figure\_x\_scale =} \DecValTok{5}\NormalTok{,}
                             \AttributeTok{Taxa =}\NormalTok{ OSLO.infant.taxa,}
                             \AttributeTok{plot\_lm =} \ConstantTok{FALSE}\NormalTok{,}
                             \AttributeTok{legend\_title\_size =} \DecValTok{10}\NormalTok{,}
                             \AttributeTok{legend\_text\_size =} \DecValTok{8}\NormalTok{)}

\CommentTok{\# Visualize the first OTU cluster in ID10}
\NormalTok{oslo\_model\_vis}\SpecialCharTok{$}\NormalTok{ID\_10[[}\DecValTok{1}\NormalTok{]]}
\end{Highlighting}
\end{Shaded}

The output results above represent OTU clusters with highly similar
temporal patterns. If you want to find other OTUs with temporal patterns
that are opposite to those of the identified OTUs, you can use the
Data.opp.cor.vis function. This function will help identify and
visualize OTUs that exhibit contrasting time patterns. The microbial
features being compared will be represented by solid curves, while those
that are negatively correlated (with a high correlation distance, for
example, greater than 0.9, and \texttt{ng\_cor\_thres} set to 0.1) will
be represented by dashed curves. This distinction will help to clearly
visualize the relationships between the microbial features in the plots.

\begin{Shaded}
\begin{Highlighting}[]
\NormalTok{opposite\_vis }\OtherTok{=} \FunctionTok{Data.opp.cor.vis}\NormalTok{(}\AttributeTok{predicted\_data =}\NormalTok{ oslo\_model\_pred,}
                                \AttributeTok{pre\_processed\_data =}\NormalTok{ oslo\_data\_trans,}
                                \AttributeTok{Design\_data =}\NormalTok{ oslo\_data\_design,}
                                \AttributeTok{ng\_cor\_thres =} \FloatTok{0.1}\NormalTok{,}
                                \AttributeTok{Taxa =}\NormalTok{ OSLO.infant.taxa,}
                                \AttributeTok{plot\_dots =} \ConstantTok{FALSE}\NormalTok{,}\AttributeTok{legend\_text\_size =} \DecValTok{5}\NormalTok{)}

\CommentTok{\# For example, taking OTU\_000001 of ID10}
\CommentTok{\# we can observe microbial features with a correlation distance greater than 0.9}
\NormalTok{opposite\_vis}\SpecialCharTok{$}\NormalTok{ID\_10}\SpecialCharTok{$}\NormalTok{OTU\_000001}
\end{Highlighting}
\end{Shaded}

\subsubsection{Case study 2: Gut microbiome dataset of preterm infant
with
sepsis.}\label{case-study-2-gut-microbiome-dataset-of-preterm-infant-with-sepsis.}

In this case study, we applied MicrobTiSDA to a subset of gut microbiome
time-series data from preterm infants diagnosed with sepsis, aiming to
investigate the dynamic differences between the gut microbiomes of
septic preterm infants infected with pathogenic E.coli and their matched
controls. This dataset consisted of 4 preterm infants diagnosed with
sepsis within the first month of life (S\_D, ID\_7, ID\_14, ID\_19,
ID\_21) and 4 individually matched controls (S\_C, ID\_31, ID\_37,
ID\_42, ID\_43).

This dataset contains a total of 51 samples and 202 OTUs, with the
metadata including information such as subject ID, sample ID, sampling
time, and sample grouping. To standardize sampling times, we normalized
the sepsis diagnosis day for each patient to the pseudo-Day 11, aligning
all sampling times to a 10-day window prior to diagnosis (i.e.,
pseudo-Days 1--10).

Load data

\begin{Shaded}
\begin{Highlighting}[]
\FunctionTok{data}\NormalTok{(}\StringTok{"preterm.data"}\NormalTok{)}
\FunctionTok{data}\NormalTok{(}\StringTok{"preterm.meta"}\NormalTok{)}
\FunctionTok{data}\NormalTok{(}\StringTok{"preterm.taxa"}\NormalTok{)}
\end{Highlighting}
\end{Shaded}

Since this case study involves comparing the dynamic differences in gut
microbiomes between two groups, we will first use the random forest
method, treating all OTUs as input features to classify the gut
microbiome samples of the two groups. We will select OTUs that
contribute significantly to the sample classification (i.e., biomarkers)
for subsequent analysis. In this case, this step will serve as a
replacement for the data filtering step.

\begin{Shaded}
\begin{Highlighting}[]
\NormalTok{preterm\_rf }\OtherTok{=} \FunctionTok{Data.rf.classifier}\NormalTok{(}\AttributeTok{raw\_data =}\NormalTok{ preterm.data,}
                                \AttributeTok{metadata =}\NormalTok{ preterm.meta,}
                                \AttributeTok{train\_p =} \FloatTok{0.7}\NormalTok{,}
                                \AttributeTok{Group =} \StringTok{\textquotesingle{}Group\textquotesingle{}}\NormalTok{,}
                                \AttributeTok{OTU\_counts\_filter\_value =} \DecValTok{0}\NormalTok{)}
\NormalTok{preterm\_rf}\SpecialCharTok{$}\NormalTok{Margin\_scores\_train}
\end{Highlighting}
\end{Shaded}

Here, we choose to use 70\% of the samples as the training set and the
remaining 30\% as the testing set. Given that the number of microbial
features in this dataset is relatively small, we didn't filter out
low\_abundant microbial features in this step
(\texttt{OTU\_counts\_filter\_value\ =\ 0}). The results indicate that
the sample classification model built using all OTUs as inputs
demonstrates high performance, with an Out of Bag error of 0.081.

Additionally, based on the cross-validation curve output above, we
selected the top 10 OTUs as biomarkers according to their mean decrease
accuracy (MDA) ranking by running the function \code{Rf.biomarkers}.

\begin{Shaded}
\begin{Highlighting}[]
\NormalTok{selected\_biomarkers }\OtherTok{=} \FunctionTok{Rf.biomarkers}\NormalTok{(}\AttributeTok{rf =}\NormalTok{ preterm\_rf,}\AttributeTok{feature\_select\_num =} \DecValTok{10}\NormalTok{)}
\end{Highlighting}
\end{Shaded}

Using the selected biomarkers and applying NMDS (Non-metric
Multidimensional Scaling), we can observe that the selected biomarkers
effectively distinguish between the two groups of samples.

\begin{Shaded}
\begin{Highlighting}[]
\NormalTok{preterm\_rf\_vis }\OtherTok{=} \FunctionTok{Classify.vis}\NormalTok{(}\AttributeTok{classified\_results =}\NormalTok{ selected\_biomarkers,}\AttributeTok{dist\_method =} \StringTok{\textquotesingle{}bray\textquotesingle{}}\NormalTok{)}
\end{Highlighting}
\end{Shaded}

Extract the biomarker data (\texttt{Important\_OTU\_table}) from the
output of the \texttt{Data.rf.classifier} function for subsequent
analysis, and first apply the modified centered log-ratio (mCLR)
transformation to the data.

\begin{Shaded}
\begin{Highlighting}[]
\NormalTok{preterm\_trans }\OtherTok{=} \FunctionTok{Data.trans}\NormalTok{(}\AttributeTok{Data =}\NormalTok{ selected\_biomarkers}\SpecialCharTok{$}\NormalTok{OTU\_importance,}
                           \AttributeTok{Group\_var =} \StringTok{\textquotesingle{}Group\textquotesingle{}}\NormalTok{,}
                           \AttributeTok{metadata =}\NormalTok{ preterm.meta)}
\end{Highlighting}
\end{Shaded}

Construct the design matrix.

\begin{Shaded}
\begin{Highlighting}[]
\NormalTok{preterm\_design }\OtherTok{=} \FunctionTok{Design}\NormalTok{(}\AttributeTok{metadata =}\NormalTok{ preterm.meta,}
                        \AttributeTok{Group\_var =} \StringTok{\textquotesingle{}Group\textquotesingle{}}\NormalTok{,}
                        \AttributeTok{Pre\_processed\_Data =}\NormalTok{ preterm\_trans,}
                        \AttributeTok{Sample\_ID =} \StringTok{\textquotesingle{}ID\textquotesingle{}}\NormalTok{,}
                        \AttributeTok{Sample\_Time =} \StringTok{\textquotesingle{}Time\textquotesingle{}}\NormalTok{)}
\end{Highlighting}
\end{Shaded}

Next, mixed-effects natural spline regression model is applied to
construct general OTU time series regression model for each group, with
time as a fixed effect and the subject ID within the group as a random
effect.

\begin{Shaded}
\begin{Highlighting}[]
\NormalTok{preterm\_model}\OtherTok{=} \FunctionTok{Reg.MESR}\NormalTok{(}\AttributeTok{Data\_for\_Reg =}\NormalTok{ preterm\_design,}
                        \AttributeTok{pre\_processed\_data =}\NormalTok{ preterm\_trans,}
                        \AttributeTok{max\_Knots =} \DecValTok{2}\NormalTok{)}
\end{Highlighting}
\end{Shaded}

Based on the fitted regression model, the \texttt{Pred.data.MESR}
function is used to predict the general time pattern of each microbial
feature within the group. Note that \texttt{Pred.data.MESR} can only
accept output from the \texttt{Reg.MESR} function.

\begin{Shaded}
\begin{Highlighting}[]
\NormalTok{preterm\_pred }\OtherTok{=} \FunctionTok{Pred.data.MESR}\NormalTok{(}\AttributeTok{Fitted\_models =}\NormalTok{ preterm\_model,}
                              \AttributeTok{metadata =}\NormalTok{ preterm.meta,}
                              \AttributeTok{Group =} \StringTok{\textquotesingle{}Group\textquotesingle{}}\NormalTok{,}
                              \AttributeTok{Sample\_Time =} \StringTok{\textquotesingle{}Time\textquotesingle{}}\NormalTok{,}
                              \AttributeTok{time\_step =} \DecValTok{1}\NormalTok{)}
\end{Highlighting}
\end{Shaded}

Following the same analytical workflow as in Case 1, we can also cluster
the predicted time patterns of microbial features.

\begin{Shaded}
\begin{Highlighting}[]
\NormalTok{preterm\_clust }\OtherTok{=} \FunctionTok{Data.cluster}\NormalTok{(}\AttributeTok{predicted\_data =}\NormalTok{ preterm\_pred,}
                             \AttributeTok{clust\_method =} \StringTok{\textquotesingle{}average\textquotesingle{}}\NormalTok{,}
                             \AttributeTok{font\_size =} \FloatTok{3.5}\NormalTok{,}
                             \AttributeTok{dend\_title\_size =} \DecValTok{12}\NormalTok{)}
\end{Highlighting}
\end{Shaded}

And we can determine the OTU clusters according to the output results
above using \code{Data.cluster.cut} by setting the cutting distance
equals to 0.2.

\begin{Shaded}
\begin{Highlighting}[]
\NormalTok{preterm\_clust\_results }\OtherTok{=} \FunctionTok{Data.cluster.cut}\NormalTok{(}\AttributeTok{cluster\_outputs =}\NormalTok{ preterm\_clust,}\AttributeTok{cut\_height =} \FloatTok{0.2}\NormalTok{,}\AttributeTok{font\_size =} \FloatTok{3.5}\NormalTok{)}
\end{Highlighting}
\end{Shaded}

Here, OTUs with a correlation distance of less than 0.2 are grouped into
clusters.

The visualization of microbial feature abundance time patterns fitted by
the linear mixed-effects spline regression model requires the use of the
\texttt{Data.visual.MESR} function.

\begin{Shaded}
\begin{Highlighting}[]
\NormalTok{preterm\_data\_vis }\OtherTok{=} \FunctionTok{Data.visual.MESR}\NormalTok{(}\AttributeTok{Design\_data =}\NormalTok{ preterm\_design,}
                                    \AttributeTok{cluster\_results =}\NormalTok{ preterm\_clust\_results,}
                                    \AttributeTok{cutree\_by =} \StringTok{\textquotesingle{}height\textquotesingle{}}\NormalTok{,}
                                    \AttributeTok{cluster\_height =} \FunctionTok{c}\NormalTok{(}\FloatTok{0.2}\NormalTok{,}\FloatTok{0.2}\NormalTok{),}
                                    \AttributeTok{predicted\_data =}\NormalTok{ preterm\_pred,}
                                    \AttributeTok{pre\_processed\_data =}\NormalTok{ preterm\_trans,}
                                    \AttributeTok{Taxa =}\NormalTok{ preterm.taxa,}
                                    \AttributeTok{plot\_dots =} \ConstantTok{TRUE}\NormalTok{,}
                                    \AttributeTok{plot\_lm =} \ConstantTok{FALSE}\NormalTok{,}
                                    \AttributeTok{figure\_x\_scale =} \DecValTok{1}\NormalTok{,}
                                    \AttributeTok{legend\_title\_size =} \DecValTok{10}\NormalTok{,}
                                    \AttributeTok{legend\_text\_size =} \DecValTok{8}\NormalTok{,}
                                    \AttributeTok{axis\_title\_size =} \DecValTok{10}\NormalTok{,}
                                    \AttributeTok{axis\_x\_size =} \DecValTok{7}\NormalTok{,}
                                    \AttributeTok{axis\_y\_size =} \DecValTok{7}\NormalTok{)}

\CommentTok{\# Example: visualize the first OTU cluster in both the S\_C and S\_D groups}
\NormalTok{preterm\_data\_vis}\SpecialCharTok{$}\NormalTok{S\_C[[}\DecValTok{1}\NormalTok{]]}
\NormalTok{preterm\_data\_vis}\SpecialCharTok{$}\NormalTok{S\_D[[}\DecValTok{1}\NormalTok{]]}
\end{Highlighting}
\end{Shaded}


\end{document}
